\uuid{Mxpr}
\titre{DL usuels/Calcul de limites/Asymptotes}
\auteur{hanani ; mzari}
\organisation{exo7}
\licence{CC BY-NC-SA 4.0}
\datecreate{2020-01-20}
\niveau{L1}
\module{Analyse}
\chapitre{Développement limité}
\sousChapitre{Applications}
\difficulte{3}
\theme{Applications des DL, DL usuels, Calcul de limites, Asymptotes}
\exo7qcmid{lille-2.developpements-limites.045}

\contenu{
\texte{
Soit $a\in \Rr$ et $f$ une fonction. On dit que $f(x)$ est un $o\left((x-a)^n\right)$ au voisinage de $a$ et on écrit $f(x)=o((x-a)^n)$ si $\displaystyle \lim _{x\to a}\frac{f(x)}{(x-a)^n}=0$. On note par $DL_n(a)f(x)$ le développement limité de $f$ à l'ordre $n$ en $a$. On note aussi $DL_n(a^+)f(x)$ (resp. $DL_n(a^-)f(x)$) le développement limité de $f$ à l'ordre $n$ à droite en $a$ (resp. à gauche en $a$).
}

\begin{question}
\qclassification{125.03}
\qtags{motcle=DL usuels/Calcul de limites/Asymptotes}
On considère les fonctions $f$ et $g$ telles que $\displaystyle f(x)=x^2\arctan \frac{1}{1+x^2}$ et $\displaystyle g(x)=f\left(\frac{1}{x}\right)$. On note $\Gamma$ le graphe de $f$ et $\Delta$ la droite d'équation $\displaystyle y=1$. Parmi les affirmations suivantes, cocher celles qui sont vraies :
\begin{answers}
\good{$\displaystyle DL_2(0^+)g(x)=1+2x+\frac{x^2}{2}+o(x^2)$.}
\bad{$\displaystyle \lim _{x\to 0^+}g(x)=1$ et $\displaystyle \lim _{x\to 0^-}g(x)=-1$.}
\bad{Au voisinage de $+\infty$, $\Gamma$ admet la droite $\Delta$ comme asymptote et $\Gamma$ est situé au dessus de $\Delta$.}
\good{Au voisinage de $-\infty$, $\Gamma$ admet la droite $\Delta$ comme asymptote et $\Gamma$ est situé en dessous de $\Delta$.}
\end{answers}
\vskip2mm
\begin{explanations}
On a : $\displaystyle g(x)=\frac{1}{x^2}\arctan \frac{x^2}{1+x^2}$. On vérifie que $\displaystyle DL_4(0)\arctan u=u-\frac{u^3}{3}+o(u^4)$ et donc avec $\displaystyle u=\frac{x^2}{1+x^2}=x^2-x^4+o(x^4)$, on obtient :
$$DL_2(0)g(x)=1-x^2+o(x^2) \Rightarrow DL_2(\pm\infty)f(x)=1-\frac{1}{x^2}+o\left(\frac{1}{x^2}\right).$$
Ainsi, au voisinage de $\pm\infty$, $\Gamma$ admet $\Delta$ comme asymptote et $\Gamma$ est situé en dessous de $\Delta$ car $\displaystyle f(x)-y=-\frac{1}{x^2}+o\left(\frac{1}{x^2}\right)\leq 0$. Enfin, $\displaystyle \lim _{x\to 0^+}g(x)=\lim _{x\to 0^-}g(x)=1$.
\end{explanations}
\end{question}
}
